% Generated draft: requires mathematical and source review.
\documentclass[11pt]{article}
\usepackage[margin=25mm]{geometry}
\usepackage{fontspec}
\setmainfont{Latin Modern Roman}
\usepackage{amsmath,amssymb,mathtools,unicode-math}
\setmathfont{Latin Modern Math}
\usepackage{enumitem,needspace,xurl,hyperref}
\setlength{\parindent}{0pt}
\setlength{\parskip}{5pt}
\setlength{\emergencystretch}{3em}
\newcommand{\N}{\mathbb{N}}
\newcommand{\Z}{\mathbb{Z}}
\newcommand{\Q}{\mathbb{Q}}
\newcommand{\R}{\mathbb{R}}
\newcommand{\C}{\mathbb{C}}
\newcommand{\F}{\mathbb{F}}
\newcommand{\A}{\mathbb{A}}
\newcommand{\PP}{\mathbb{P}}
\newcommand{\E}{\mathbb{E}}
\newcommand{\eps}{\varepsilon}
\newcommand{\dd}{\,\mathrm{d}}
\begin{document}
\section*{UnsolvedMath ID: 30004580}
\textbf{Source problem code:} catalogue entry
\subsection{Definitions and mathematical statement}
The supplied record poses a historical research question, rather than a mathematical conjecture with specified objects, hypotheses, and quantifiers. Its subject is the demographics of practitioners in medieval mathematical professions across different contexts, together with the effects of those demographics on mathematical work, collaboration, and communication.

The source question is:
\begin{quote}
In medieval mathematical professions across different contexts, what were practitioners' demographics, and how did those demographics affect their work, collaboration, and communication?
\end{quote}

The record does not define the chronological or geographical scope of ``medieval,'' the contexts or professions to be included, the demographic attributes to be studied, or the evidence and criteria for identifying demographic effects. Consequently, a precise mathematical proposition cannot be extracted without adding unsupported assumptions. The question is retained as a historical research target; its comparative scope does not specify a universal quantifier over all medieval contexts.
\subsection{Short English statement}
Investigate the demographics of medieval mathematical practitioners across different contexts and how those demographics affected their work, collaboration, and communication. The supplied record leaves the historical scope, demographic categories, and criteria for assessing effects unspecified.
\subsection{Sources}
\begin{itemize}
\item Local UnsolvedMath registry, \texttt{data/unsolvedmath/problems.json}: permanent ID \(30004580\), title ``Demographics of Medieval Mathematical Practitioners,'' source problem code \texttt{catalogue entry}, collection ``Oberwolfach Reports Open Problems.'' The record labels the statement exact and internally complete; these are supplied registry annotations, not independent verification against the report. No individual proposer is supplied.
\item Supplied source citation: \emph{History of Mathematics: A Global Cultural Approach} (2021), Oberwolfach Reports, DOI 10.4171/owr/2020/41, pp.\ 2021--2092. Supplied source URL and literature source labeled ``original Oberwolfach report'': \url{https://doi.org/10.4171/owr/2020/41}. The background describes Oberwolfach Reports as records of mathematical workshops and their open-problem discussions. The registry gives proposed year 2020; the supplied citation gives 2021.
\item Supplied literature assessment, checked on \texttt{2026-08-22}: status ``open.'' The assessment characterizes this as a question about historical demographics and their effects across mathematical practice, recorded by the report as an open research question rather than a settled proposition. This is supplied dated curation derived from the OWR record, not independently verified here. The supplied curation requests independent verification of time-sensitive status and cited claims.
\end{itemize}
\end{document}
